% page 230 of the facsimile (image p-229 of the scan)
% block 160.0 x 246.3 mm ; left margin 8.0 mm
\pgc{-12.700mm}{0.000mm}{
\l{\cel{2}222.}
\l{}
\l{}
\l{\sou{homilio}. I. (\sou{liturgio katolika}). Instruktado familiara di la}
\l{\cel{3}populo pri evangelio e pri la materii religiala. - II.}
\l{\cel{3}Breviario-leciono extraktita de la homilii di la patri}
\l{\cel{3}ekleziala. - III. (desprizoze). Etiko-prediko, leciono}
\l{\cel{3}enoyiganta. - DEFIS.}
\l{}
\l{homocentro.\cel{2}(matem.) Centro komuna di plura centri.}
\l{}
\l{homogena. Di qua la elementi konstitucanta esas sama-natura.}
\l{\cel{3}- DEFIRS.}
\l{}
\l{\sou{homografio}. (\sou{geom.)}\cel{2}Interdependo partikulara di du figuri}
\l{\cel{3}geometriala.}
\l{}
\l{\sou{homologio}. (\sou{geom.)}\cel{2}Karaktero di la figuri tala (segun Ponce-}
\l{\cel{3}let) ke la punti interkorespondanta di ica e di ita esas}
\l{\cel{3}duopa sur la rekti qui kunkuras ad un punto (centro di homo-}
\l{\cel{3}logio) e ke la rekti qui juntas du punti ek ica e la du}
\l{\cel{3}korespondanti di la altra iras krucumar sur rekto una (axo}
\l{\cel{3}di homologio) se la figuri esas plana, o sur un plano (plano}
\l{\cel{3}di homologio) se la figuri esas irgequala en spaco.}
\l{}
\l{\sou{homologa}.\cel{2}(\sou{aludante un o plura termini}). Qui interkorespondas.}
\l{\cel{3}DEFIRS.}
\l{}
\l{\sou{homonima}. (\sou{gram.}) Vorti qui esas intersama ortografiale ed}
\l{\cel{3}pronuncale, ma interdiferas pri la senco. -\cel{1}\sou{Homonimo}\cel{1}: per-}
\l{\cel{3}sono di qua la nomo esas sama kam olta di altra persono. -}
\l{\cel{3}DEFIRS.}
\l{}
\l{\sou{homoteta}. (\sou{geom.}) Donite un punto fixa O\cel{2}ed un konstanto k,}
\l{\cel{3}se, ad un punto M ek la spaco onu korespondantigas un punto}
\l{\cel{3}M\textquotesingle{}, tala ke\cel{1}\sou{OM\textquotesingle{}/OM}\cel{1}egalesas k, onu dicas ke la punti M e M\textquotesingle{}}
\l{\cel{3}esas \textquotedbl{}homoteta\textquotedbl{} pri la punto O.}
\l{}
\l{\sou{homotipa}. (\sou{anat.}) Di qua la organi, en un individuo, esas\cel{2}la}
\l{\cel{3}analogi di altra organi. (kom ex.: la pedo-fingri kompare}
\l{\cel{3}kun la manuo-fingri).}
\l{}
\l{\sou{honesta}. Qua esas konforma a la devo, a vertuo. - EFIS.}
\l{}
\l{\sou{honoro}. I. Digneso etikala qua naskas de la bezono di estimeso}
\l{\cel{3}da la kunhomi e da ni ipsa. - II. Distingo-manifesto por}
\l{\cel{3}ulu, flatema. - EFIS.}
\l{}
\l{\sou{honorario}. Pago-pekunio po la servi da persono qua havas pro-}
\l{\cel{3}fesiono honorizanta (advokato, mediko, sacerdoto, e c.)}
\l{\cel{3}- DEFIRS.}
\l{}
\l{\sou{hop}\cel{1}!\cel{2}Interjeciono uzata por stimular, ecitar a kuresko, o}
\l{\cel{3}saltigar trans obstakli.}
\l{}
\l{\sou{hoplito}. (en la epoki\cel{1}\sou{antiqua)}\cel{2}Infantriano di qua la armi esis}
\l{\cel{3}tre pezoza. - DEFIRS.}
}
